El modelo BPMN de la Figura~\ref{fig:asis_subproceso_2_bpmn} representa el mismo proceso distribuido
en tres carriles (\emph{swimlanes}): Comisi\'on de Matr\'icula, Sistema Intranet DIRCE /
SIGA-ORCE, y Estudiante Regular. Esta distribuci\'on permite identificar con claridad las
responsabilidades de cada actor y los puntos de interacci\'on entre el estudiante y el
sistema institucional.

\begin{figure}[H]
    \centering
    \includegraphics[width=0.95\textwidth]{asis_subproceso_2_bpmn.png}
    \caption{Modelo BPMN AS IS -- Subproceso 2: Selecci\'on y Registro de Cursos.}
    \label{fig:asis_subproceso_2_bpmn}
\end{figure}

En el carril de la \textbf{Comisi\'on de Matr\'icula} se configura previamente la oferta
acad\'emica, los turnos por facultad y los cupos por secci\'on. Esta actividad precede al
inicio formal del subproceso. El carril del \textbf{Sistema (Intranet DIRCE / SIGA-ORCE)}
concentra las actividades automatizadas: autenticar al estudiante, desplegar la oferta,
ejecutar las validaciones de cr\'editos, horarios y vacantes, y generar la Pre-Ficha de
matr\'icula. El carril del \textbf{Estudiante Regular} comprende las actividades
interactivas: ingresar al sistema, seleccionar y ajustar cursos, y confirmar la
inscripci\'on.

%% ----------------------------------------------------------
